% !TeX root = part3.tex

\documentclass[../final.tex]{subfiles}

\begin{document}

% При отдельной сборке продолжаем нумерацию полного конспекта.
\ifSubfilesClassLoaded{\setcounter{chapter}{1}\setcounter{section}{0}}{}


% ============================================================
% Квантовая электродинамика
% Лекция 08.10
% ============================================================

\chapter{Квантовая электродинамика}

\rsubtitle{Скалярная электродинамика}

Основная литература к этой главе ---
Р.~Фейнман, \emph{<<Теория фундаментальных процессов>>},
а также Ф.~Халзен, А.~Мартин, \emph{<<Кварки и лептоны>>}
(обе книги лежат в папке <<Литература>>).
Сначала рассматриваются частицы без спина (уравнение Клейна--Гордона),
на них отрабатывается вся техника: теория возмущений, амплитуды,
сечения, диаграммы. Частицы со спином вводятся в конце главы.


% ============================================================
\section{Нерелятивистская квантовая механика}
% ============================================================

% ------------------------------------------------------------
\subsection{Уравнение Шрёдингера}
% ------------------------------------------------------------

Состояние частицы описывается волновой функцией, подчиняющейся
\rconcept{уравнению Шрёдингера}

\begin{equation*}
    i\hbar\frac{\partial\psi}{\partial t}
    =
    \hat H\psi.
\end{equation*}

Далее везде используется \rassumption{система единиц $\hbar=c=1$}.
Для свободной частицы $\hat H=\hat{\mathbf p}^2/2m$,
$\hat{\mathbf p}=-i\nabla$, и уравнение принимает вид

\begin{equation*}
    i\frac{\partial\psi}{\partial t}
    +
    \frac{1}{2m}\nabla^2\psi
    =
    0.
\end{equation*}

Величина $|\psi|^2=\rho$ имеет смысл плотности вероятности.


% ------------------------------------------------------------
\subsection{Уравнение непрерывности}
% ------------------------------------------------------------

Умножая уравнение на $\psi^*$, сопряжённое уравнение на $\psi$
и вычитая, получаем \rconcept{уравнение непрерывности}

\begin{equation*}
    \frac{\partial\rho}{\partial t}
    +
    \nabla\cdot\mathbf j
    =
    0,
\end{equation*}

где плотность потока вероятности

\begin{equation*}
    \mathbf j
    =
    -\frac{i}{2m}
    \left(
        \psi^*\nabla\psi
        -
        \psi\nabla\psi^*
    \right).
\end{equation*}

Решение для свободной частицы --- плоская волна

\begin{equation*}
    \psi(\mathbf r,t)
    =
    N e^{i\mathbf p\mathbf r-iEt},
\end{equation*}

для которой $\rho=|N|^2$, $\mathbf j=|N|^2\,\mathbf p/m$.

\remark
Уравнение Шрёдингера \rnote{не является лоренц-инвариантным}:
время входит в него в первой степени производной, а координаты ---
во второй. Для релятивистских частиц нужно другое уравнение.


% ============================================================
\section{Уравнение Клейна--Гордона}
% ============================================================

% ------------------------------------------------------------
\subsection{Вывод уравнения}
% ------------------------------------------------------------

Исходим из релятивистской связи энергии и импульса

\begin{equation*}
    E^2-\mathbf p^2=m^2
\end{equation*}

и заменяем величины операторами

\begin{equation*}
    \hat E=i\frac{\partial}{\partial t},
    \qquad
    \hat{\mathbf p}=-i\nabla.
\end{equation*}

Получаем

\begin{equation*}
    -\frac{\partial^2\psi}{\partial t^2}
    +
    \nabla^2\psi
    =
    m^2\psi
\end{equation*}

--- \rconcept{уравнение Клейна--Гордона}. В ковариантной записи

\begin{equation*}
    \left(\Box+m^2\right)\varphi=0,
    \qquad
    \Box
    =
    \partial_\mu\partial^\mu
    =
    \frac{\partial^2}{\partial t^2}-\nabla^2,
\end{equation*}

где

\begin{equation*}
    \partial^\mu
    =
    \frac{\partial}{\partial x_\mu}
    =
    \left(
        \frac{\partial}{\partial t},\,
        -\nabla
    \right).
\end{equation*}


% ------------------------------------------------------------
\subsection{Плотность и ток}
% ------------------------------------------------------------

Аналогично нерелятивистскому случаю из уравнения следует

\begin{equation*}
    \frac{\partial}{\partial t}
    \underbrace{
        \left[
            i\left(
                \psi^*\frac{\partial\psi}{\partial t}
                -
                \psi\frac{\partial\psi^*}{\partial t}
            \right)
        \right]
    }_{\rho}
    +
    \nabla\cdot
    \underbrace{
        \left[
            -i\left(
                \psi^*\nabla\psi
                -
                \psi\nabla\psi^*
            \right)
        \right]
    }_{\mathbf j}
    =
    0.
\end{equation*}

В четырёхмерной форме

\begin{equation*}
    j^\mu
    =
    i\left(
        \varphi^*\partial^\mu\varphi
        -
        \varphi\,\partial^\mu\varphi^*
    \right),
    \qquad
    \partial_\mu j^\mu=0.
\end{equation*}

Решение для свободной частицы ---

\begin{equation*}
    \varphi=Ne^{-ipx},
    \qquad
    px=Et-\mathbf p\mathbf x.
\end{equation*}

Подставляя, находим

\begin{equation*}
    j^\mu=2|N|^2p^\mu,
    \qquad
    \rho=2E|N|^2.
\end{equation*}

\remark
Плотность $\rho$ пропорциональна энергии, а уравнение
допускает решения с $E=\pm\sqrt{\mathbf p^2+m^2}$.
Поэтому \rnote{$\rho$ не является положительно определённой}
и не может трактоваться как плотность вероятности.
Решения с $E<0$ интерпретируются как античастицы
(раздел~\ref{sec:antiparticles}), а $e\rho$ --- как плотность
заряда, которая может иметь любой знак.

\end{document}
